\chapter{Findings and Lessons Learned}
\label{ch:lessons}

\section{Summary of findings}
\label{sec:findings}

The research question asked whether a local language model places better chunk boundaries
than rule-based chunking techniques, and under which conditions. The evaluation of
\cref{ch:evaluation} answers it.

\paragraph{Better than splitters that cut sentences apart.} The LLM chunker retrieves
significantly better than fixed-size splitting and LangChain's default recursive splitter:
on \code{rfc9110} by $+16.4$~pp at \hitk{1} against fixed-size splitting, and on
\code{nasa} by $+8.2$ and $+7.9$~pp against the two. These gains mainly stem from the fact
that both baselines cut inside sentences. On \code{wells}, whose short paragraphs the
default recursive splitter keeps whole, no difference is significant.

\paragraph{Not better than splitters that keep sentences intact.} Against the recursive
splitter with sentence ends, the tuned semantic chunker and the same chunker without the
model, no significant difference remains. The LLM chunker leads by 0.4 to 4.7~pp on
\code{nasa} and \code{rfc9110} and is behind in five of six comparisons on \code{wells},
by up to 5.8~pp. None of these differences is significant, which indicates that once a
baseline keeps sentences intact, the decisions of the local model bring no measurable
gain.

\paragraph{The optional features do not improve retrieval.} The low-information filter
lowers \hitk{3} by 3.6~pp on \code{nasa} and 1.8~pp on \code{rfc9110} (\cref{sec:filter}).
Heading detection changes no hit rate by more than 0.7~pp, topic labels lower \hitk{1} on
all three documents, and letting the model choose the split point at the size cap makes no
significant difference (\cref{sec:ablations}). Parent-child retrieval raises most hit
rates, but the rule-based splitters gain more from it than the LLM chunker
(\cref{sec:parentchild}).

\paragraph{Cost and conclusion.} Although run times depend strongly on the hardware, LLM
chunking takes orders of magnitude longer than rule-based splitting: hours instead of about
a second per document. For retrieval as measured here, a length-matched recursive splitter
that prefers sentence ends to spaces reaches comparable hit rates, and combined with
parent-child retrieval it reaches a higher \hitk{1} than the LLM chunker on all three
documents. The cost of LLM-based chunking is paid only once, at indexing time, since
afterwards the chunks only have to be retrieved from the vector store. On the benchmarks
performed here, however, the difference in retrieval was not large enough to justify that
overhead.

\section{Lessons learned}
\label{sec:lessons}

The most challenging part of the project was the fair evaluation: building baselines that
are really comparable, measuring significance and separating the contribution of the
language model from that of the sentence splitting. The following lessons come from that
part.

\subsection{A fair comparison needs more than equal chunk length}

At first, only the mean chunk length of the baselines was adjusted to the LLM chunker.
Evaluating the first results revealed gaps in the experimental design: LangChain's
semantic chunker had to be tuned before it was a fair opponent, the recursive splitter
needed a variant that prefers sentence ends, and the same chunker without the model had to
be added to isolate the contribution of the model.

\subsection{Changes to the plan will happen}

The test configuration changed over time. For example, exact search replaced the
approximate index, and more baselines and a family of 30 tests were added after an
external review. Furthermore, the low-information filter was removed from the tested arm
because it removed passages that contained answers.

\subsection{Cached results ensure clean tracking}

Chunking takes hours, so the chunks have to be cached, and each cache needs a clear
fingerprint of its configuration and code, so that results from different states are not
mixed up.

\subsection{Working with limited hardware}

All runs used a laptop with \SI{8}{\giga\byte} of memory, and chunking one document took
two to six hours. Computing the embeddings in Ollama instead of the Python process and
caching every arm separately made the runs possible. The run time is one of the reasons
why the evaluation covers only three documents.

\section{Limitations and future work}
\label{sec:future}

The evaluation covers one document per structure class, the chunking model is a local
model with four billion parameters, and the embedding model and the parameter settings
were fixed for the entire project. Furthermore, the benchmark measures whether a verbatim
passage is retrieved intact, with questions generated by a language model; it does not
measure the quality of generated answers or whether a chunk reads as one coherent unit.
The results on \code{nasa} depend on an extraction that joins each page into one
paragraph, and the ablations were run only with the low-information filter.

Future work can address these limitations with several documents per structure class, a
larger chunking model, an end-to-end evaluation of the generated answers, a repetition on
\code{nasa} with a layout-preserving extraction and ablations without the filter. For the
library, returning the offsets of each chunk and the reason for each boundary (topic
decision, size cap or heading) instead of plain strings would make its decisions easier to
inspect in an application.
